\documentclass[11pt,a4paper]{article}
\usepackage[T1]{fontenc}
\usepackage{lmodern}
\usepackage[a4paper,margin=2.5cm]{geometry}
\usepackage{microtype}
\usepackage{amsmath,amssymb,amsthm,mathtools}
\usepackage{booktabs}
\usepackage{enumitem}
\usepackage[hidelinks]{hyperref}
\usepackage[capitalise,noabbrev]{cleveref}

\newtheorem{theorem}{Theorem}[section]
\newtheorem{lemma}[theorem]{Lemma}
\newtheorem{proposition}[theorem]{Proposition}
\newtheorem{corollary}[theorem]{Corollary}
\theoremstyle{definition}
\newtheorem{definition}[theorem]{Definition}
\newtheorem{example}[theorem]{Example}
\newtheorem{assumption}[theorem]{Assumption}
\theoremstyle{remark}
\newtheorem{remark}[theorem]{Remark}

\newcommand{\N}{\mathbb{N}}
\newcommand{\Z}{\mathbb{Z}}
\newcommand{\Q}{\mathbb{Q}}
\newcommand{\E}{\mathbb{E}}
\renewcommand{\P}{\mathbb{P}}
\newcommand{\pX}{\ensuremath{\mathord{+}\mathsf{x}}}
\newcommand{\mX}{\ensuremath{\mathord{-}\mathsf{x}}}
\newcommand{\pU}{\ensuremath{\mathord{+}\mathsf{1}}}
\newcommand{\mU}{\ensuremath{\mathord{-}\mathsf{1}}}
\newcommand{\dbl}{\ensuremath{\mathord{\times}\mathsf{2}}}
\newcommand{\hlv}{\ensuremath{\mathord{\div}\mathsf{2}}}
\newcommand{\sq}{\ensuremath{\surd}}
\newcommand{\abs}[1]{\lvert #1\rvert}
\newcommand{\norm}[1]{\lVert #1\rVert}
\newcommand{\pos}[1]{{#1}^{+}}
\DeclareMathOperator{\sgn}{sgn}

\title{The economy of \emph{Yhtälösade}:\\[2pt]
\large when can a resource-limited equation game never be lost?}
\author{Technical note for the Matikkapelit project}
\date{17 October 2026}

\begin{document}
\maketitle

\begin{abstract}
\noindent
In \emph{Yhtälösade} linear equations fall down a well and the player solves them by
dragging operation bricks ($\pX$, $\mX$, $\pU$, $\mU$) onto them; every use costs one unit
of a finite stock. A solved equation $x=\xi$ becomes a bomb: for $\xi>0$ it returns $\xi$
units to the stock, split at random between the bricks, and for $\xi<0$ it clears trash.
We ask which constraints on the solutions $\xi_1,\dots,\xi_N$ and on the operations
needed make the game impossible to lose for a player who solves optimally.
We show that for the current level the cheapest solution is unique, so the game reduces
to an exact ledger of demand and reward vectors.
A conservation law then yields a balance inequality that any never-losing design must
satisfy. The current parameters violate it by $457/404\approx1.13$ units per
equation, and even a clairvoyant reward rule could not save more than $16\%$ of
20-equation games. The random split makes things worse: we prove that under it no
sure guarantee exists short of a starting stock covering all future demand.
We then give two designs with sure guarantees, both with proofs:
a \emph{deficit-filling} reward rule, whose condition is a Lindley-queue inequality on
the sequence $(\delta_i-\pos{\xi_i})$, and an \emph{adaptive generator} that works
with any reward rule.
Finally we extend the cost theory to doubling, halving and square roots. The cost of
an equation then becomes a capped signed-binary weight, the order of the operations
starts to matter, and $\pm$-bombs act as automatic stabilisers of the economy.
\end{abstract}

\tableofcontents

%======================================================================
\section{Introduction}\label{sec:intro}

\emph{Yhtälösade} (``equation rain'') is a Tetris-like game for Finnish upper-secondary
students. One equation at a time falls from the top of a well with brick walls. Four
special bricks in the walls carry the operations $\pX$, $\mX$, $\pU$ and $\mU$, each with
a count. Dragging a brick onto the falling equation applies its operation to both sides
and lowers the count by one. When the equation reads $x=\xi$ it turns into a bomb
labelled $\xi$:
\begin{itemize}[nosep]
\item a positive bomb returns $\xi$ units of operations, split at random between the bricks;
\item a negative bomb clears trash around the place where it lands.
\end{itemize}
An equation that reaches the ground unsolved stays there as trash, and trash above a red
line ends the game.

In the first test game the designer, who believes they played optimally, ran out of
$\mX$ at the third equation and out of $\mU$ at the fifth. After that only $\pU$ was
left, and almost nothing could be solved. This note explains why that outcome is
typical, not bad luck. It then answers the designer's question: \emph{what constraints
on the solutions $\xi_1,\dots,\xi_N$ and on the operations they require guarantee that
an optimally playing student never loses?}

\paragraph{Results in brief.}
\begin{enumerate}[label=(\roman*),nosep]
\item On the current level the optimal play for each equation is unique up to order
  and uses a demand vector $d(E)$ that does not depend on the solution
  (\cref{prop:cost}). The game therefore reduces to an exact ledger
  (\cref{sec:ledger}).
\item \emph{Balance law} (\cref{thm:balance}). A perfect player survives $N$ rounds
  only if, for every $i$,
  $\sum_{t\le i}\delta_t-\sum_{t<i}\pos{\xi_t}\le\norm{c_0}$, where $\delta_t$ is the
  number of operations equation $t$ needs and $\norm{c_0}$ the starting stock. In the
  long run this means $(1-q)\,\nu_+\ge\bar\delta$, where $q$ is the share of negative
  solutions and $\nu_+$ the mean positive solution. The current generator has
  $(1-q)\nu_+=1.75$ against $\bar\delta=2.88$.
\item \emph{Random rewards cannot be relied on} (\cref{thm:diffuse}). Under the current
  random split, survival is certain only if the starting stock covers the total future
  demand, brick by brick. A perfect player solves all 20 equations in only about $0.2\%$
  of games, and the first stall comes on average at equation $5.5$
  (\cref{tab:sim}).
\item \emph{Sure guarantees.} With a deficit-filling reward rule,
  survival is guaranteed whenever a Lindley queue driven by $(\delta_i-\pos{\xi_i})$
  stays below the starting buffer (\cref{thm:main}). One concrete instance: every
  negative solution is followed by a positive one that repays it. With $c_0=(5,5,5,5)$
  this keeps half of the solutions negative and all of them in $\pm\{1,\dots,6\}$
  (\cref{cor:alternating}). Separately, a generator that looks at the stock guarantees
  survival under \emph{any} reward rule (\cref{thm:adaptive}).
\item \emph{Extensions} (\cref{sec:double,sec:halve,sec:sqrt}). With $\dbl$ or $\hlv$ the
  cheapest cost is $m+w_m(A)+w_m(B)$, with $w_m$ a capped signed-binary weight. The order
  of the operations matters, and the cheapest demand vector need not be unique.
  Square roots create $\pm$-bombs that both refuel and clean: a third of symmetric
  quadratics is enough to cancel the current deficit.
\end{enumerate}
All numerical claims were checked by exact computation, by exhaustive search or by
Monte Carlo simulation; see \cref{app:methods}.

%======================================================================
\section{The model}\label{sec:model}

\subsection{Equations and operations}
Throughout, an \emph{equation} is $U(x)=L(x)$ with affine sides and
\begin{equation}\label{eq:diff}
U(x)-L(x)=\Delta\,(x-\xi),\qquad \Delta>0,\ \xi\in\Z\setminus\{0\}.
\end{equation}
$U$ is the side with the larger coefficient of $x$ and $\xi$ is the solution. (To avoid a
clash with the variable $x$, we write $\xi_i$ for the solutions $x_1,\dots,x_N$ of the
question.) On the current level $\Delta=1$. The \emph{bricks} form a finite set $J$, on
the current level $J=\{\pX,\mX,\pU,\mU\}$, and brick $j$ applies its operation to both
sides. A \emph{stock} (inventory) is a vector $c\in\N^J$, and its \emph{total} is
$\norm{c}:=\sum_j c(j)$, the number of operations it holds. This is the $\ell^1$ norm,
which on $\N^J$ is the plain sum; we use it for all vectors in $\N^J$ (stocks, demands,
rewards, counts of operations used), while $\abs{z}$ is always the absolute value of an
integer.
Inequalities between vectors are componentwise, and $\pos z=\max(z,0)$,
$z^-=\max(-z,0)$.

\subsection{A round of the game}
Equations $E_1,E_2,\dots$ fall one at a time; the next one appears only after the
previous one has been resolved. In round $i$, with stock $c_{i-1}$:
\begin{itemize}[nosep]
\item each operation the player applies costs one unit of its brick, and only a brick with
  a positive count can be used;
\item as soon as the equation reads $x=\xi$ or $\xi=x$ it locks and becomes a bomb;
\item a bomb with $\xi>0$ adds a \emph{reward} $r_i\in\N^J$ with $\norm{r_i}=\xi$ to the
  stock, chosen by an \emph{allocation rule};
\item a bomb with $\xi<0$ adds nothing but clears trash (\cref{sec:trash});
\item an unsolved equation that lands becomes a trash block one row high.
\end{itemize}

\paragraph{The current level (``D0'').} The parameters are:
\begin{itemize}[nosep]
\item starting stock $c_0=(3,3,5,5)$ (order $\pX,\mX,\pU,\mU$) and $N=20$ rounds;
\item $\xi$ uniform on $\pm\{1,\dots,6\}$;
\item each equation is built backwards from $x=\xi$ by adding $k_b$ boxes and a
  constant $k$ to both sides: $(1+k_b)x+k=k_bx+\xi+k$, with the sides swapped at random;
\item $k_b$ uniform on the multiset $\{-2,-1,-1,0,1,1,2\}$, $k$ uniform on
  $\{-3,\dots,3\}$; pairs with $\abs{k_b}+\abs{k}\notin[1,4]$ are rejected, and those with
  $\abs{k_b}+\abs{k}=1$ are kept with probability $2/5$;
\item the random split rule: $k\sim U\{1,\dots,\min(\xi,4)\}$ parts, each starting at 1,
  the remaining $\xi-k$ units added to uniformly random parts, and the parts sent to $k$
  distinct random bricks.
\end{itemize}

\begin{definition}[Perfect player, stall, survival]\label{def:play}
The \emph{perfect player} has no time pressure, wastes no operations, and solves the
current equation whenever some solving sequence is affordable, using a cheapest
affordable one. Round $i$ \emph{stalls} if no solving sequence is affordable with
$c_{i-1}$. The game is \emph{strongly survived} over $N$ rounds if no round stalls.
\end{definition}

``Never lost with optimal bombing'' is taken to mean strong survival of the perfect
player. \Cref{sec:trash} shows that the game's formal end, Game Over, is out of reach in
a 20-equation game on a typical phone. The loss players actually experience is the
stall: once the stock is depleted, every later equation turns into trash.

%======================================================================
\section{Solving cost on the current level}\label{sec:cost}

\begin{lemma}[Invariants]\label{lem:inv}
Every operation preserves $U-L$, hence $\Delta$ and $\xi$, and $U$ stays the side with the
larger coefficient.
\end{lemma}
\begin{proof}
Each of $\pX,\mX,\pU,\mU$ adds the same affine function $g\in\{x,-x,1,-1\}$ to both sides,
so $(U+g)-(L+g)=U-L$.
\end{proof}

By \eqref{eq:diff} with $\Delta=1$, every level-one equation with solution $\xi$ can be
written uniquely as
\begin{equation}\label{eq:pq}
U=(1+p)x+q,\qquad L=px+(\xi+q),\qquad (p,q)\in\Z^2 .
\end{equation}
The generator produces $(p,q)=(k_b,k)$.

\begin{lemma}[Coordinates]\label{lem:coord}
The operations act on $(p,q)$ as unit translations:
$\pX\colon p\mapsto p+1$, $\mX\colon p\mapsto p-1$, $\pU\colon q\mapsto q+1$ and
$\mU\colon q\mapsto q-1$. The equation reads $x=\xi$ or $\xi=x$ if and only if
$(p,q)=(0,0)$.
\end{lemma}
\begin{proof}
The first claim is immediate from \eqref{eq:pq}. For the second, a solved equation has
one side equal to $x$ and the other a constant. By \cref{lem:inv} the coefficient of $U$
exceeds that of $L$ by one, so $U=x$ and $L=\xi$, i.e.\ $1+p=1$ and $q=0$. Conversely
$(p,q)=(0,0)$ gives $U=x$, $L=\xi$.
\end{proof}

\begin{proposition}[Cost and forced demand]\label{prop:cost}
Let $n_j$ be the number of uses of brick $j$ in a sequence of operations. The sequence
solves $E=(p,q;\xi)$ if and only if $n_{\mX}-n_{\pX}=p$ and $n_{\mU}-n_{\pU}=q$.
Consequently
\[
n\;\ge\;d(E):=(p^-,\;\pos p,\;q^-,\;\pos q)
\]
componentwise for every solving sequence. The cheapest cost is
$\delta(E):=\norm{d(E)}=\abs p+\abs q$, and it is attained exactly by the sequences
that use $d(E)$, in any order.
\end{proposition}
\begin{proof}
The translations commute, so the final point is
$(p-n_{\mX}+n_{\pX},\,q-n_{\mU}+n_{\pU})$, and \cref{lem:coord} gives the first claim.
From $n_{\mX}=p+n_{\pX}$ we get $n_{\mX}\ge p$ and $n_{\mX}\ge0$, so
$n_{\mX}\ge\pos p$; likewise $n_{\pX}\ge p^-$, $n_{\mU}\ge\pos q$ and $n_{\pU}\ge q^-$.
Summing gives $\norm n\ge\abs p+\abs q$, with equality if and only if every inequality is
tight, i.e.\ $n=d(E)$. A minimal sequence never passes through $(0,0)$ before its end,
so the lock-on-solve rule does not interfere.
\end{proof}

\begin{corollary}[Optimal bombing is unique and dominant]\label{cor:dominant}
On the current level the perfect player solves $E$ with exactly $d(E)$. Among all ways of
solving $E$ this leaves the largest stock in every brick. Hence the whole economy is
described by the pairs $(d_i,\xi_i)$, $d_i:=d(E_i)$, $\delta_i:=\norm{d_i}$.
\end{corollary}

\begin{remark}[Negative solutions cost the same]\label{rem:symmetry}
The demand $d(E)$ does not depend on $\xi$. The reflection
$(p,q,\xi)\mapsto(p,-q,-\xi)$, which turns $ax+b=cx+d$ into $ax-b=cx-d$, preserves
$\delta$ and only swaps the roles of $\pU$ and $\mU$.
\end{remark}

\begin{table}[t]\centering\small
\caption{Exact demand statistics of the generators.
D0 is the current level; D1 is the rebalanced proposal of \cref{sec:recommend}
($\abs{k_b}\le2$, $\abs k\le 2$, $\delta\le3$, cost-1 pairs kept with probability $1/2$,
$\xi\in\{-6,\dots,-1\}\cup\{3,\dots,8\}$).}\label{tab:demand}
\begin{tabular}{lcccccc}
\toprule
 & $\P(\delta{=}1,2,3,4)$ & $\E\delta$ & $\E d(\pm\mathsf x)$ & $\E d(\pm\mathsf 1)$ & $\E\pos\xi$ & $\E[\pos\xi-\delta]$\\
\midrule
D0 & $\tfrac{6}{101},\tfrac{30}{101},\tfrac{35}{101},\tfrac{30}{101}$ & $\tfrac{291}{101}\approx2.881$ & $\tfrac{57}{101}\approx0.564$ & $\tfrac{177}{202}\approx0.876$ & $\tfrac74$ & $-\tfrac{457}{404}\approx-1.131$\\[2pt]
D1 & $\tfrac19,\tfrac49,\tfrac49,0$ & $\tfrac73\approx2.333$ & $\tfrac59\approx0.556$ & $\tfrac{11}{18}\approx0.611$ & $\tfrac{11}4$ & $+\tfrac5{12}\approx+0.417$\\
\bottomrule
\end{tabular}
\end{table}

%======================================================================
\section{The ledger and the balance law}\label{sec:ledger}

By \cref{cor:dominant}, as long as no round stalls the stock evolves as
\begin{equation}\label{eq:ledger}
c_i=c_{i-1}-d_i+r_i,\qquad \norm{r_i}=\pos{\xi_i},
\end{equation}
and round $i$ stalls if and only if $c_{i-1}\not\ge d_i$.

\begin{proposition}[Conservation]\label{prop:conservation}
If rounds $1,\dots,i$ are all solved by the perfect player, then
$\norm{c_i}=\norm{c_0}+\sum_{t\le i}(\pos{\xi_t}-\delta_t)$. A real player who wastes
$W_i\ge0$ operations ends with $W_i$ units less.
\end{proposition}
\begin{proof}
Take the norm of \eqref{eq:ledger}, using $\norm{d_t}=\delta_t$ and $\norm{r_t}=\pos{\xi_t}$.
\end{proof}

\begin{theorem}[Balance law]\label{thm:balance}
\leavevmode
\begin{enumerate}[label=(\alph*),nosep]
\item If the perfect player strongly survives rounds $1,\dots,N$, then for every $i\le N$
  \begin{equation}\label{eq:balance}
  \sum_{t\le i}\delta_t-\sum_{t<i}\pos{\xi_t}\;\le\;\norm{c_0}.
  \end{equation}
\item If an infinite sequence is strongly survived, then
  $\liminf_N\frac1N\sum_{t\le N}(\pos{\xi_t}-\delta_t)\ge0$. In terms of the empirical
  share $q$ of negative solutions, the mean positive solution $\nu_+$ and the mean
  demand $\bar\delta$, this is the long-run condition
  \begin{equation}\label{eq:balance-mean}
  (1-q)\,\nu_+\;\ge\;\bar\delta .
  \end{equation}
  If in addition the stock stays bounded, \eqref{eq:balance-mean} holds with equality in
  the limit.
\end{enumerate}
\end{theorem}
\begin{proof}
(a) Round $i$ is solvable only if $c_{i-1}\ge d_i$, which forces
$\norm{c_{i-1}}\ge\delta_i$; insert \cref{prop:conservation}.
(b) Dividing \eqref{eq:balance} by $N$ gives
$\frac1N\sum_{t<N}(\pos{\xi_t}-\delta_t)\ge(\delta_N-\norm{c_0})/N\ge-\norm{c_0}/N$.
Moreover $\frac1N\sum\pos{\xi_t}=(1-q)\nu_+$ and $\frac1N\sum\delta_t=\bar\delta$.
A bounded stock also bounds the sum from above by \cref{prop:conservation}.
\end{proof}

The balance law is a property of the sequence alone, so no reward rule can repair a
violation of it. For the current level, \cref{tab:demand} gives
$\E[\pos\xi-\delta]=-457/404$. By the strong law of large numbers, every long game
therefore stalls almost surely. For $N=20$ the bound is quantitative:

\begin{theorem}[Allocation-free upper bound]\label{thm:allocfree}
Let the rounds be i.i.d. Under \emph{any} allocation rule, even one that knows the
future,
\[
\P(\text{strong survival over }N\text{ rounds})\;\le\;
\bar p_N:=\P\Bigl(\forall i\le N:\ \norm{c_0}+\sum_{t<i}(\pos{\xi_t}-\delta_t)\ge\delta_i\Bigr).
\]
For D0 with $\norm{c_0}=16$, exact dynamic programming over the total stock gives
$\bar p_5=0.987$, $\bar p_{10}=0.634$ and $\bar p_{20}=0.162$.
\end{theorem}
\begin{proof}
By \cref{thm:balance}(a), strong survival implies the event in $\bar p_N$, and that event
depends only on $(\delta_t,\xi_t)$, not on how rewards are split.
\end{proof}

So with the current generator at most $16\%$ of 20-equation games could be survived even
by an ideal reward rule; the random split achieves $0.17\%$ (\cref{tab:sim}). The gap is
due to misallocation, which is the subject of \cref{sec:random}. Once the stock is low,
the game cannot recover by itself:

\begin{proposition}[Poverty trap]\label{prop:trap}
The empty stock is absorbing, since every equation has $\delta\ge1$. Under D0, a stock
holding only $\pU$ (at least three units) solves the next equation with probability
$6/101\approx0.059$, and earns a reward with probability $3/101\approx0.030$.
\end{proposition}
\begin{proof}
With only $\pU$ available, $d(E)\le c$ means $p=0$ and $-3\le q<0$. From the
generator's joint law (\cref{app:methods}) these pairs have total probability
$\frac1{101}+\frac5{202}+\frac5{202}=\frac6{101}$, and $\xi>0$ independently with
probability $\frac12$.
\end{proof}

This is exactly the designer's end state: about $34$ rounds per rewarded solve on
average.

%======================================================================
\section{Random allocation}\label{sec:random}

\begin{definition}
An allocation rule is \emph{$\eta$-diffuse} ($\eta>0$) if, for every positive bomb and
every brick $j$, $\P(r(j)=0\mid\text{past})\ge\eta$.
\end{definition}

\begin{lemma}\label{lem:diffuse}
The current split rule, implemented with an unbiased shuffle, is $\tfrac38$-diffuse and
gives $\E[r(j)\mid\xi]=\xi/4$ for every $j$.
\end{lemma}
\begin{proof}
With $M=\min(\xi,4)$ and $k$ parts on $k$ random bricks out of four,
$\P(r(j)=0)=\frac1M\sum_{k=1}^M(1-\frac k4)=1-\frac{M+1}8\ge\frac38$.
The bricks are exchangeable, so each brick receives $\xi/4$ on average.
\end{proof}

\begin{theorem}[Random rewards cannot be relied on]\label{thm:diffuse}
Fix a sequence $E_1,\dots,E_N$ and an $\eta$-diffuse rule. Strong survival has
probability one if and only if $c_0\ge\sum_{t\le N}d_t$.
\end{theorem}
\begin{proof}
($\Leftarrow$) Rewards are non-negative, so as long as no round has stalled,
$c_{i-1}\ge c_0-\sum_{t<i}d_t\ge d_i$.
($\Rightarrow$) Suppose $c_0(j)<\sum_{t\le N}d_t(j)$ and let $i^*$ be the first $i$
with $\sum_{t\le i}d_t(j)>c_0(j)$. Let $A$ be the event that no bomb before round $i^*$
sends anything to $j$. By the chain rule and diffuseness, $\P(A)\ge\eta^{i^*-1}>0$. On
$A$, either an earlier round stalls, or
$c_{i^*-1}(j)=c_0(j)-\sum_{t<i^*}d_t(j)<d_{i^*}(j)$ and round $i^*$ stalls.
\end{proof}

Under D0 the expected 20-round demand on $\pU$ alone is $20\cdot\frac{177}{202}\approx17.5$,
against $c_0(\pU)=5$, so a sure guarantee is out of the question with random rewards.
The best one can hope for is a probabilistic guarantee, and its sign is decided brick by
brick by the \emph{drift}
\[
\mu_j:=\E r(j)-\E d(j)\qquad(\text{per round, }r(j):=0\text{ when }\xi\le0).
\]

\begin{theorem}[Negative drift: exponential collapse]\label{thm:negdrift}
Let the rounds be i.i.d. and let rewards depend only on $\xi$ and fresh randomness. If
$\mu_j<0$ for some $j$, then for $N>c_0(j)/\abs{\mu_j}$
\[
\P(\text{strong survival over }N)\le
\exp\Bigl(-\frac{2\,(N\abs{\mu_j}-c_0(j))^2}{N\,s_j^2}\Bigr),
\]
where $s_j$ is the range of $r(j)-d(j)$.
\end{theorem}
\begin{proof}
Let $X_t=r_t(j)-d_t(j)$ (the reward drawn as if round $t$ were solved). The $X_t$ are
i.i.d.\ with mean $\mu_j$. On the survival event \eqref{eq:ledger} gives
$0\le c_N(j)=c_0(j)+\sum_{t\le N}X_t$, so
$\sum_t(X_t-\mu_j)\ge N\abs{\mu_j}-c_0(j)>0$. Hoeffding's inequality bounds the
probability of this.
\end{proof}

For D0 with the unbiased split, $\mu_{\pm\mathsf x}=\frac7{16}-\frac{57}{101}=-\frac{205}{1616}\approx-0.127$
and $\mu_{\pm\mathsf 1}=\frac7{16}-\frac{177}{202}=-\frac{709}{1616}\approx-0.439$: every brick
drifts down, and the constant bricks drift fastest. Hoeffding is loose at $N=20$;
simulation (\cref{tab:sim}) shows that only $0.17\%$ of games are survived and that the
first stall comes at round $5.5$ on average. When all drifts are positive, the classical
Cramér--Lundberg argument gives a bound that holds for an unlimited number of rounds:

\begin{theorem}[Positive drift: Lundberg bound]\label{thm:lundberg}
Under the assumptions of \cref{thm:negdrift}, let $\theta_j>0$ satisfy
$\E[e^{-\theta_j r(j)}]\,\E[e^{\theta_j d(j)}]\le1$. Then the probability that brick $j$
ever causes a stall is at most $e^{-\theta_j(c_0(j)-d_j^{\max}+1)}$, and
$\P(\text{some stall ever})\le\sum_j e^{-\theta_j(c_0(j)-d_j^{\max}+1)}$.
\end{theorem}
\begin{proof}
Let $W_i=c_0(j)+\sum_{t<i}r_t(j)-\sum_{t\le i}d_t(j)$. If the first stall happens in
round $i$ because of $j$, then all earlier rounds were solved and $W_i=c_{i-1}(j)-d_i(j)<0$.
Write $W_i=W_1+\sum_{t=2}^iZ_t$ with $Z_t=r_{t-1}(j)-d_t(j)$. Different $Z_t$ involve
disjoint rounds, and within a round $\xi$ and $(k_b,k)$ are drawn independently, so the
$Z_t$ are i.i.d.\ with $\E e^{-\theta_jZ}\le1$. Hence $M_i=e^{-\theta_j(W_i-W_1)}$ is a
non-negative supermartingale. Let $\tau=\inf\{i:W_i<0\}$; since $W$ is integer-valued,
$M_\tau\ge e^{\theta_j(W_1+1)}$ on $\{\tau<\infty\}$. Optional stopping gives
$\P(\tau\le n)\,e^{\theta_j(W_1+1)}\le\E M_{\tau\wedge n}\le1$. Let $n\to\infty$ and use
$W_1\ge c_0(j)-d_j^{\max}$. The last claim is a union bound.
\end{proof}

\begin{example}
For the rebalanced generator D1 with the random split, all drifts are positive
($\frac{19}{144}$ for $\pm\mathsf x$ and $\frac{11}{144}$ for $\pm\mathsf 1$). The adjustment
coefficients are $\theta_{\pm\mathsf x}=0.127$ and $\theta_{\pm\mathsf1}=0.066$. To push the
Lundberg bound below $0.35$, a stock of $30$ \emph{per brick} is needed.
Probabilistic guarantees with random rewards are expensive, because the buffer grows like
$\theta^{-1}\log(K/\varepsilon)$.
\end{example}

\begin{remark}[Matching principle]\label{rem:matching}
If a rule sends a share $\pi_j$ of each reward to brick $j$, then
$\mu_j=\pi_j\E\pos\xi-\E d(j)$. The uniform shares $\pi_j=\frac14$ ignore that, in D0,
the constant bricks are needed $1.55$ times as often as the $x$-bricks. Making every
$\mu_j\ge0$ then needs $\E\pos\xi\ge4\max_j\E d(j)=3.50$, whereas the balance law only
asks for $2.88$. Demand-proportional shares $\pi_j=\E d(j)/\E\delta$ give
$\mu_j=\pi_j(\E\pos\xi-\E\delta)$, so every brick inherits the sign of the global
balance.
\end{remark}

\begin{remark}[Implementation detail]\label{rem:shuffle}
The game shuffles the bricks by sorting them with a random comparator,
\[\texttt{sort(() => Math.random() - .5)},\]
which is not a uniform shuffle. Measured in V8 (Node 22, the engine of Chrome), a one-part reward goes
to $\pX$, $\mX$, $\pU$, $\mU$ with probabilities $0.36$, $0.14$, $0.19$, $0.31$. This
lowers $\mu_{\mX}$ from $-0.127$ to about $-0.160$. In the test game $\mX$ was the first
brick to run out. A Fisher--Yates shuffle removes the bias.
\end{remark}

%======================================================================
\section{Designs with sure guarantees}\label{sec:sure}

\Cref{thm:diffuse} shows that a sure guarantee needs either a reward rule that knows
what is needed, or a generator that knows what is in stock. We treat both.

\subsection{Deficit filling and the Lindley queue}

\begin{definition}
\emph{Deficit filling} sends each reward unit, one at a time, to a brick minimising
$c(j)-c_0(j)$, the count relative to its starting level (ties broken arbitrarily). The
\emph{total deficit} is $T(c)=\sum_j(c_0(j)-c(j))^+$.
\end{definition}

\begin{lemma}\label{lem:lindley}
Under deficit filling, if round $i$ is solved then
$T(c_i)\le\bigl(T(c_{i-1})+\delta_i-\pos{\xi_i}\bigr)^+$.
\end{lemma}
\begin{proof}
Paying $d_i$ raises each term $(c_0(j)-c(j))^+$ by at most $d_i(j)$, so $T$ grows by at
most $\delta_i$. Now add the $\pos{\xi_i}$ reward units one by one. While $T>0$, some brick
is below its starting level, so the chosen brick is too, and the unit lowers $T$ by
exactly one. When $T=0$ it stays $0$.
\end{proof}

\begin{theorem}[Main theorem: sure survival under deficit filling]\label{thm:main}
Define the Lindley sequence $\Lambda_0=0$,
$\Lambda_i=(\Lambda_{i-1}+\delta_i-\pos{\xi_i})^+$. Under deficit filling the perfect
player strongly survives rounds $1,\dots,N$ whenever
\begin{equation}\label{eq:main}
\Lambda_{i-1}\;\le\;\min_{j}\bigl(c_0(j)-d_i(j)\bigr)\qquad(1\le i\le N).
\end{equation}
In particular it suffices that $\min_jc_0(j)\ge L+\max_{i<N}\Lambda_i$, where
$L=\max_{i,j}d_i(j)$.
\end{theorem}
\begin{proof}
By induction, assume rounds $1,\dots,i-1$ were solved. Since $z\mapsto(z+a)^+$ is
monotone, \cref{lem:lindley} gives $T(c_{i-1})\le\Lambda_{i-1}$. Then for every $j$,
$c_{i-1}(j)\ge c_0(j)-(c_0(j)-c_{i-1}(j))^+\ge c_0(j)-T(c_{i-1})\ge c_0(j)-\Lambda_{i-1}\ge d_i(j)$
by \eqref{eq:main}, so round $i$ is solvable.
\end{proof}

\Cref{eq:main} is the requested constraint on the solutions: it couples each
$\xi_i$ to the demands only through the scalar queue $\Lambda$. A negative $\xi_i$ is a
\emph{debt} $\delta_i$ joining the queue, and a positive $\xi_i$ with
$\xi_i>\delta_i$ \emph{serves} it. The simplest sequences keep the queue short by placing
each debt directly before its repayment:

\begin{corollary}[Debt, then repayment]\label{cor:alternating}
Suppose the rounds form consecutive blocks of two kinds:
\begin{itemize}[nosep]
\item a single round with $\xi\ge\delta$;
\item a pair (negative round, positive round) whose positive solution repays both,
  $\xi_{\rm pos}\ge\delta_{\rm pos}+\delta_{\rm neg}$.
\end{itemize}
Then $\Lambda_i\le P:=\max\delta$ for all $i$, so $c_0(j)\ge L+P$ guarantees strong
survival under deficit filling. With $P=3$ and $\abs{k_b},\abs k\le2$ (so $L=2$),
$c_0=(5,5,5,5)$ is enough. All solutions stay in $\pm\{1,\dots,6\}$, because a
repayment needs at most $\delta_{\rm pos}+\delta_{\rm neg}\le6$, and strict alternation
makes exactly half of the solutions negative.
\end{corollary}
\begin{proof}
Single blocks keep $\Lambda=0$. In a pair, $\Lambda=\delta_{\rm neg}\le P$ after the
negative round and $\Lambda=(\delta_{\rm neg}+\delta_{\rm pos}-\xi_{\rm pos})^+=0$ after
the positive one. Condition \eqref{eq:main} reads $0\le c_0(j)-d_i(j)$ in the first
round of a block and $\delta_{\rm neg}\le c_0(j)-d_i(j)$ in the second, both implied
by $c_0(j)\ge L+P$.
\end{proof}

\begin{remark}[Repay, do not prepay]
Before any debt exists, deficit filling spreads surplus evenly over the $K$ bricks, so
raising every brick by $D$ costs $KD$ units. Repaying an existing debt of $D$ costs $D$.
Debts should therefore come \emph{before} their repayments. This is convenient: the
generator, which builds equations backwards, already knows the debt when it picks the
next positive solution.
\end{remark}

\begin{proposition}[Random sequences under deficit filling]
If the rounds are i.i.d.\ with $\E\pos\xi>\E\delta$, then $\Lambda$ is a stable queue and
$\P(\Lambda_i\ge b)\le e^{-\theta b}$ for all $i$, where $\theta>0$ solves
$\E e^{\theta(\delta-\pos\xi)}=1$. Hence
$\P(\text{a stall within }N\text{ rounds})\le N\,e^{-\theta(\min_jc_0(j)-L+1)}$.
\end{proposition}
\begin{proof}
By the Lindley representation, $\Lambda_i$ has the law of
$\max_{0\le k\le i}$ of the partial sums of the reversed increments, which is dominated
by the supremum of a random walk with negative drift. The Lundberg argument of
\cref{thm:lundberg} bounds that supremum. A stall in round $i$ needs
$\Lambda_{i-1}\ge\min_jc_0(j)-L+1$; take a union bound over $i\le N$.
\end{proof}

For D1, $\theta=0.092$; the bound is loose, but simulation shows a large real effect
(\cref{tab:sim}). With $c_0=(4,4,4,4)$, deficit filling raises the share of survived
20-round games from $16\%$ to $66\%$, and with $c_0=(6,6,6,6)$ to $89\%$.

\subsection{Adaptive generation}

\begin{definition}
Fix $a\ge1$. The \emph{adaptive generator} $G_a$ draws round $i$ from a base generator
conditioned on the admissible set
\[
\mathcal A(c_{i-1})=\bigl\{(p,q,\xi):\ d(p,q)\le c_{i-1},\ \ \pos\xi\ge\delta(p,q)+a-\norm{c_{i-1}}\bigr\}.
\]
\end{definition}

\begin{theorem}[Sure survival under any allocation]\label{thm:adaptive}
Suppose the base generator gives positive probability to every cost-one pair and to
some positive solution (independently of the pair), and that $\norm{c_0}\ge a$. Then $G_a$ is well defined in every round, and the
perfect player strongly survives every number of rounds, under \emph{every} allocation
rule, adversarial ones included.
\end{theorem}
\begin{proof}
By induction, assume $\norm{c_{i-1}}\ge a\ge1$. Some brick $j$ has $c_{i-1}(j)\ge1$, and
the cost-one pair using only $j$, together with any positive $\xi$ of the base law,
satisfies $\pos\xi\ge1\ge1+a-\norm{c_{i-1}}$. So $\mathcal A(c_{i-1})$ has positive mass and round
$i$ is drawn from it. Its equation is affordable because $d_i\le c_{i-1}$. Whatever the
split, $\norm{r_i}=\pos{\xi_i}$, so
$\norm{c_i}=\norm{c_{i-1}}-\delta_i+\pos{\xi_i}\ge a$.
\end{proof}

The guarantee is unconditional, but by \cref{thm:balance} the generator can only buy it
by skewing the sequence: under a deficit it must offer more positive solutions or
cheaper equations. Simulation measures this skew (\cref{tab:sim}). On D0, $G_1$ lowers the
share of negative solutions from $0.50$ to $0.47$ (20 rounds) and $0.44$ (100 rounds),
and the mean demand from $2.88$ to $2.34$ and $2.09$. On the balanced D1, $G_4$ is
almost invisible: $q=0.49$--$0.50$, $\bar\delta=2.25$--$2.27$. The adaptive generator is
best used as a safety net under a generator that already satisfies the balance law.

%======================================================================
\section{Trash, Game Over and the value of a negative bomb}\label{sec:trash}

Let the well have $R$ rows, with Game Over when a trash block lands in one of the top
two rows (above the red line). Write $H=R-2$.

\begin{lemma}[Height]\label{lem:height}
Each trash landing raises the stack height by at most one, while clearing and settling
never raise it. Game Over therefore needs at least $H+1$ trash blocks.
\end{lemma}
\begin{proof}
A block lands one row above the highest column it covers. So if the stack height is $h$,
it lands in row $r\ge R-h-1$ and the new height is at most $h+1$. Clearing removes
bricks and settling moves them down. Game Over means $r<2$, which requires
$h\ge R-2=H$ before the landing, i.e.\ at least $H$ earlier landings.
\end{proof}

The game sizes the well to the screen. Measured values of $H$ are 18 (390×844),
19 (412×915), 16 (360×740), 13 (375×667) and 12 (1366×860). Equation blocks are
$3,4,5$ cells wide with frequencies $0.38,0.57,0.05$, and on 9 columns they overlap
heavily: the simulated mean number of trash blocks to Game Over is $23.6$ for $H=18$
and $16.1$ for $H=12$. Under D0 with $H=18$, no simulated 20-round game reached Game
Over, even counting no clearing at all. With $H=12$ the rate is $5.5\%$.

\begin{proposition}[Clearing depth]\label{prop:clear}
A negative bomb $\xi$ clears the trash bricks whose centres lie within
$\rho(\xi)=0.8+0.7\abs\xi$ cells of the bomb's centre. If it lands on a trash column at
least $\lfloor\rho(\xi)\rfloor$ deep, that column loses its top $\lfloor\rho(\xi)\rfloor$ bricks, and in deep,
flat trash it removes $\#\{(i,j)\in\Z\times\Z_{\ge1}:i^2+j^2\le\rho^2\}$ bricks:
\begin{center}\small
\begin{tabular}{lcccccc}\toprule
$\abs\xi$&1&2&3&4&5&6\\\midrule
column depth&1&2&2&3&4&4\rlap{\,\textsuperscript{$*$}}\\
bricks in flat trash&3&4&10&15&26&30\\\bottomrule
\end{tabular}
\end{center}
\textsuperscript{$*$}In floating point $0.8+0.7\cdot6=4.999\ldots$, so the depth is 4, not 5.
\end{proposition}
\begin{proof}
The bomb's bottom rests on the column surface $s$ and its centre is $\frac12$ cell
higher. A brick in the same column at row $s-1+j$ ($j\ge1$) is at distance $j$, and one
$i$ columns aside at distance $\sqrt{i^2+j^2}$. Count the lattice points.
\end{proof}

\begin{proposition}[Paradox of the negative bomb]\label{prop:paradox}
In any design that guarantees strong survival, the perfect player never produces trash,
so their negative bombs never clear anything. Each negative round is then a pure cost of
$\delta_i>0$ units.
\end{proposition}

The proof is immediate. The proposition says that clearing only has value for
imperfect play: time pressure, wasted operations, slow recognition. That is exactly
where trash comes from in real games, so negative bombs act as the second chance for
human mistakes, and their cost must be built into the balance law. Under the current
rules a player can even \emph{gain} by letting negative equations fall unsolved: each
one costs a single trash block, and by \cref{lem:height} a phone with $H=18$ tolerates
18 of them. The square-root bombs of \cref{sec:sqrt} resolve this tension, because they
refuel and clean at the same time.

%======================================================================
\section{Adding doubling (\texorpdfstring{$\times2$}{x2})}\label{sec:double}

We now allow equations such as $\frac12x=5$ or $\frac12x+1=6$ and a brick $\dbl$ that
multiplies both sides by $2$. Write an equation as $U=\alpha x+\beta$ with
$L=(\alpha-\Delta)x+(\beta+\Delta\xi)$, and let $u:=\alpha-\Delta$ be the coefficient of
$x$ in $L$. The state is $s=(\alpha,\beta,\Delta)\in\Q^3$ with $\xi$ fixed. The
operations act as
\[
\pX,\mX\colon\alpha\mapsto\alpha\pm1,\qquad
\pU,\mU\colon\beta\mapsto\beta\pm1,\qquad
\dbl\colon s\mapsto2s,\qquad(\hlv\colon s\mapsto s/2\text{ in \cref{sec:halve}}),
\]
and, by the argument of \cref{lem:coord}, the equation is solved if and only if
$s=(1,0,1)$.

\begin{lemma}[Levels]\label{lem:levels}
In a sequence of operations, the \emph{level} $\ell_k$ of a translation $k$ is the number
of doublings minus the number of halvings applied after it. If the sequence contains
$\sigma$ more doublings than halvings, its final state is
\[
\Bigl(2^\sigma\alpha+\sum_{k\in X}\varepsilon_k2^{\ell_k},\;
2^\sigma\beta+\sum_{k\in C}\varepsilon_k2^{\ell_k},\;2^\sigma\Delta\Bigr),
\]
where $X$ and $C$ are the $x$- and constant translations and $\varepsilon_k=\pm1$ their
signs.
\end{lemma}
\begin{proof}
All operations are affine maps $v\mapsto\lambda v+w$ on $\Q^3$. A translation applied at
some moment is afterwards multiplied by $2$ at every doubling and by $\frac12$ at every
halving; induct over the sequence.
\end{proof}

\begin{definition}
For $m\in\N$ and $Z\in\Z$, the \emph{$m$-capped signed-binary weight} $w_m(Z)$ is the least
number of terms $\pm2^\ell$ with $0\le\ell\le m$ that sum to $Z$.
\end{definition}

\begin{lemma}[Properties of $w_m$]\label{lem:weight}
\leavevmode
\begin{enumerate}[label=(\roman*),nosep]
\item $w_0(Z)=\abs Z$. For $m\ge1$, $w_m(Z)=w_{m-1}(Z/2)$ if $Z$ is even, and
  $w_m(Z)=1+\min\{w_{m-1}(\tfrac{Z-1}2),\,w_{m-1}(\tfrac{Z+1}2)\}$ if $Z$ is odd.
\item $\lceil\abs Z/2^m\rceil\le w_m(Z)\le\lfloor\abs Z/2^m\rfloor+m$.
\item $w_m(2^mZ)=\abs Z$, and only terms $\pm2^m$ of one sign attain it.
\end{enumerate}
\end{lemma}
\begin{proof}
(i) Take an optimal representation. Two level-0 terms of opposite signs cancel and can
be deleted. Two of the same sign can be replaced by one level-1 term. So an optimum has
at most one level-0 term, and the parity of $Z$ decides whether it has one. Removing it
and halving the rest gives a representation of $(Z\mp1)/2$ or $Z/2$ with levels in
$[0,m-1]$, and conversely.
(ii) Each term is at most $2^m$ in absolute value. For the upper bound, use
$\lfloor\abs Z/2^m\rfloor$ terms $\pm2^m$ and the binary digits of the remainder, which is
below $2^m$.
(iii) The lower bound in (ii) is $\abs Z$, and it is attained only if every term has
absolute value $2^m$ and the same sign.
\end{proof}

\begin{theorem}[Cost with doubling]\label{thm:double}
With $J'=\{\pX,\mX,\pU,\mU,\dbl\}$, an equation is solvable if and only if
$\Delta=2^{-m}$ for some $m\in\N$ and $A:=2^mu$, $B:=2^m\beta$ are integers. Its
cheapest cost is
\[
C(E)=m+w_m(A)+w_m(B).
\]
\end{theorem}
\begin{proof}
Being solved requires $2^\sigma\Delta=1$ by \cref{lem:levels}, and $\sigma\ge0$ is the
number of doublings, so $\Delta=2^{-m}$ with exactly $m$ doublings and all levels in
$[0,m]$. Since $2^m\Delta=1$, the solved condition $\alpha_{\rm final}=1$ reads
$\sum_X\varepsilon_k2^{\ell_k}=1-2^m\alpha=-A$, and $\beta_{\rm final}=0$ reads
$\sum_C\varepsilon_k2^{\ell_k}=-B$. Both sums are integers, which gives necessity, and
every solving sequence costs at least $m+w_m(A)+w_m(B)$ because $w_m$ is even in $Z$.
Conversely, take optimal representations of $-A$ and $-B$. For $\ell=m,m-1,\dots,0$,
perform the translations of level $\ell$ and then (if $\ell>0$) one doubling. This
realises the bound.
\end{proof}

Both formulas, this one and the halving version of \cref{thm:halve}, agree with
breadth-first search over all $17\,474$ states within 9 operations of the solved state,
with no exceptions (\cref{app:methods}).

\begin{proposition}[Consequences of doubling]\label{prop:double-econ}
\leavevmode
\begin{enumerate}[label=(\alph*),nosep]
\item \emph{Order matters.} If $u,\beta\in\Z$ (integer constants), then
  $C(E)=m+\abs u+\abs\beta$, attained only by doing all translations \emph{before} any
  doubling. The textbook move of doubling first to clear the fraction costs
  $m+2^m(\abs u+\abs\beta)$, which is $(2^m-1)(\abs u+\abs\beta)$ extra units. For
  example $\frac12x+1=6$ takes 2 operations as $\mU,\dbl$ but 3 as $\dbl,\mU,\mU$.
\item \emph{Larger bombs.} If all displayed constants are integers, then
  $\xi\in2^m\Z$, since $\Delta\xi$ is a difference of integer constants. A doubling
  equation therefore pays $\abs\xi\ge2^m$: with solutions in $\pm\{1,\dots,6\}$ and
  $m=1$, only $\pm2,\pm4,\pm6$ occur, with mean positive reward $4$.
\item \emph{Separate budget.} The $\dbl$ brick is consumed exactly $m$ times per
  equation, so it needs its own balance: $\E r(\dbl)\ge\E m$.
\end{enumerate}
\end{proposition}
\begin{proof}
(a) By \cref{lem:weight}(iii) with $A=2^mu$ and $B=2^m\beta$, only level-$m$ terms are
optimal, i.e.\ translations made before the first doubling. Doubling first puts all
translations at level $0$, where each is worth $1$.
(b) The constant of $L$ is $\beta+\Delta\xi$; if it and $\beta$ are integers, so is
$2^{-m}\xi$.
(c) Every solving sequence contains exactly $m$ doublings.
\end{proof}

%======================================================================
\section{Adding halving (\texorpdfstring{$\div2$}{:2})}\label{sec:halve}

Halving allows integer equations with $\Delta=2^m$, such as $2x+4=10$,
$3x+1=x+7$ or $4x+3=11$.

\begin{theorem}[Cost with halving]\label{thm:halve}
With $J''=\{\pX,\mX,\pU,\mU,\hlv\}$, an equation is solvable if and only if $\Delta=2^m$
for some $m\in\N$ and $u,\beta\in\Z$. Its cheapest cost is $C(E)=m+w_m(u)+w_m(\beta)$.
\end{theorem}
\begin{proof}
Now $\sigma=-m$ and the levels lie in $[-m,0]$. The solved condition reads
$\sum_X\varepsilon_k2^{\ell_k}=1-2^{-m}\alpha=-2^{-m}u$. Multiplying by $2^m$ gives
$\sum_X\varepsilon_k2^{\ell_k+m}=-u$ with exponents in $[0,m]$, and likewise for $\beta$.
The rest is as in \cref{thm:double}. In words, a translation made after $h$ halvings is
worth $2^h$ in the original equation.
\end{proof}

\begin{example}\label{ex:halve}
$2x+4=10$: $1+w_1(4)=3$ ($\hlv,\mU,\mU$; subtracting first takes 5).
$2x+3=7$: $1+w_1(3)=3$ ($\mU,\hlv,\mU$: make the constant even first).
$3x+1=x+7$: $1+w_1(1)+w_1(1)=3$ ($\mX,\mU,\hlv$).
$4x+3=11$: $2+w_2(3)=4$, with \emph{two incomparable} cheapest demands:
$\mU,\hlv,\mU,\hlv$ uses $(\mU{:}2,\hlv{:}2)$, and $\pU,\hlv,\hlv,\mU$ uses
$(\pU{:}1,\mU{:}1,\hlv{:}2)$.
\end{example}

\begin{proposition}[Economics of halving]\label{prop:halve-econ}
\leavevmode
\begin{enumerate}[label=(\alph*),nosep]
\item Halving cuts the constant demand from $\abs\beta$ to $w_m(\beta)\in
  [\lceil\abs\beta/2^m\rceil,\lfloor\abs\beta/2^m\rfloor+m]$, at the price of $m$
  units of $\hlv$, and it puts no restriction on $\xi$. Large constants become cheap.
\item The set $\mathcal D(E)$ of demand vectors of solving sequences can have several
  incomparable minimal elements (\cref{ex:halve}). The equation is affordable if and
  only if $c\in\mathcal D(E)+\N^{J''}$, so \cref{cor:dominant} fails: the perfect player
  has to \emph{choose} which bricks to spend. In \cref{ex:halve} a stock with one $\mU$
  and one $\pU$ can only use the second solution.
\end{enumerate}
\end{proposition}
\begin{proof}
(a) is \cref{lem:weight}(ii) together with $\xi$ being free in $L$'s constant
$\beta+2^m\xi$. (b) is the example.
\end{proof}

\begin{proposition}[Both scalings: detours]\label{prop:detour}
With both $\dbl$ and $\hlv$ available, an equation with $\Delta=1$ might be solved more
cheaply by a detour through other scales. The cheapest cost is still $\abs u+\abs\beta$
whenever $\abs u+\abs\beta\le5$, and the bound is sharp: $x+6=\xi+6$ takes 5 operations
($\hlv$, three $\mU$, $\dbl$) instead of 6.
\end{proposition}
\begin{proof}
A sequence with scalings has $h\ge1$ doublings and $h$ halvings, and by
\cref{lem:levels} its translation terms are at most $2^h$ in absolute value. Its cost is
therefore at least
$g(h)=2h+\lceil\abs u/2^h\rceil+\lceil\abs\beta/2^h\rceil$. Write
$S=\abs u+\abs\beta\le5$. For $h=1$,
$g\ge2+\lceil S/2\rceil\ge S$ (check $S=5$: $2+3=5$). For $h\ge2$,
$g\ge4+[S>0]\ge S$. The example is verified directly, and the search in
\cref{app:methods} finds no shorter detour for $S\le5$.
\end{proof}

%======================================================================
\section{Adding square roots: \texorpdfstring{$\pm$}{plus-minus}-bombs}\label{sec:sqrt}

Consider equations $(x-s)^2+k=t^2+k$ with $s,k\in\Z$ and $t\in\Z_{>0}$, for example
$x^2=4$, $x^2+1=10$ or $(x-1)^2=9$. A brick $\sq$ applies only to an equation of the form
$(x-s)^2=t^2$ and turns it into $x-s=\pm t$, which the game carries on as one
$\pm$-equation: both branches receive the same operations.

\begin{lemma}[Cost of a quadratic]\label{lem:sqrtcost}
With the bricks $\pX,\mX,\pU,\mU,\sq$, the cheapest cost is $\abs k+1+\abs s$. It uses
$\abs k$ units of $\mU$ if $k>0$ (of $\pU$ if $k<0$) before the root, one $\sq$, and
$\abs s$ units of $\pU$ if $s>0$ (of $\mU$ if $s<0$) after it. Solved, it gives the two
solutions $\xi_{1,2}=s\pm t$.
\end{lemma}
\begin{proof}
Before $\sq$ only translations act. To reach the form ``square $=$ constant'', the
$x$-translations must cancel, since otherwise the right side contains $x$, and the
constant translations must total $-k$. That takes at least $\abs k$ operations, and
exactly $\abs k$ without $x$-operations. After $\sq$ both branches are level-one
equations with $(p,q)=(0,-s)$, so \cref{prop:cost} gives $\abs s$ further operations.
\end{proof}

A solved $\pm$-equation produces a \emph{$\pm$-bomb}, best shown as two bombs falling
side by side. It returns $\varrho=\pos{(s+t)}+\pos{(s-t)}$ units and clears with each
negative root.

\begin{proposition}[Self-financing quadratics]\label{prop:selffin}
A quadratic pays for itself ($\varrho\ge\abs k+1+\abs s$) if and only if:
\begin{itemize}[nosep]
\item for $\abs s<t$ (roots of opposite signs): $t\ge\abs k+1+\abs s-s$, i.e.\
  $t\ge\abs k+1$ when $s\ge0$ and $t\ge\abs k+1+2\abs s$ when $s<0$;
\item for $s>t$ (both roots positive): $s\ge\abs k+1$;
\item for $s<-t$ (both roots negative): never.
\end{itemize}
($s=\pm t$ gives the root $0$, which the generator excludes.)
In particular a symmetric quadratic ($s=0$) with $t\ge\abs k+1$ both refuels and
produces the negative bomb $-t$.
\end{proposition}
\begin{proof}
For $\abs s<t$, $\varrho=s+t$. For $s>t$, $\varrho=2s$ and $\abs s=s$. For $s<-t$,
$\varrho=0$.
\end{proof}

\begin{theorem}[Quadratics as stabilisers]\label{thm:stabilise}
Mix the linear rounds of a generator with drift $\mu_{\rm lin}=\E[\pos\xi-\delta]<0$ and
a fraction $\varphi$ of symmetric quadratics with drift
$\mu_{\rm q}=\E t-1-\E\abs k>0$. The mixture satisfies the balance law in expectation if
and only if
\[
\varphi\;\ge\;\frac{\abs{\mu_{\rm lin}}}{\abs{\mu_{\rm lin}}+\mu_{\rm q}}.
\]
For D0 ($\mu_{\rm lin}=-\frac{457}{404}$) with $t$ uniform on $\{2,\dots,6\}$ and $k$
uniform on $\{-1,0,1\}$ ($\mu_{\rm q}=\frac73$), the threshold is $\varphi\ge0.327$.
Every quadratic round also yields one negative bomb, so the share of negative bombs
stays at one half.
\end{theorem}
\begin{proof}
By linearity the mixed drift is $(1-\varphi)\mu_{\rm lin}+\varphi\mu_{\rm q}$; solve
for $\varphi\ge0$. The numbers follow from \cref{tab:demand}.
\end{proof}

Square roots therefore resolve the paradox of \cref{prop:paradox}: the negative half of a
$\pm$-bomb is paid for by its positive half. The $\sq$ brick needs its own budget
($\E r(\sq)\ge\varphi$). The game must refuse $\sq$ on a non-square, such as
$x^2+5=9$ before the 5 is removed, and that refusal is itself a lesson.

%======================================================================
\section{Recommendations}\label{sec:recommend}

\begin{table}[t]\centering\small
\caption{Monte Carlo results for a perfect player (20{,}000 runs for $N=20$, 4{,}000
for $N=100$, 5{,}000 for D2). ``Survival'' is the share of games without a stall. DF is
deficit filling, $G_a$ the adaptive generator and D2 the debt-then-repay design of
\cref{cor:alternating}. With the biased shuffle of the actual implementation, D0
survives $0.14\%$ of 20-round games (200{,}000 runs).}\label{tab:sim}
\begin{tabular}{lccccccc}
\toprule
& & \multicolumn{3}{c}{$N=20$} & \multicolumn{2}{c}{$N=100$} & \\
\cmidrule(lr){3-5}\cmidrule(lr){6-7}
design & $c_0$ & survival & solved & 1st stall & survival & solved & $q$, $\bar\delta$\\
\midrule
D0 random split      & 3,3,5,5 & 0.002 & 9.6  & 5.5  & 0.000 & 18.3 & 0.50, 2.88\\
D0 + DF              & 3,3,5,5 & 0.035 & 11.3 & 7.5  & 0.000 & 17.2 & 0.50, 2.88\\
D0 + $G_1$           & 3,3,5,5 & 1     & 20   & --   & 1     & 100  & 0.47, 2.34 \;/\; 0.44, 2.09\\
D1 random split      & 3,3,5,5 & 0.156 & 15.0 & 7.9  & 0.028 & 67.5 & 0.50, 2.33\\
D1 random split      & 4,4,4,4 & 0.162 & 15.2 & 8.4  & 0.031 & 68.1 & 0.50, 2.33\\
D1 + DF              & 4,4,4,4 & 0.656 & 18.4 & 9.2  & 0.518 & 86.7 & 0.50, 2.33\\
D1 + DF              & 6,6,6,6 & 0.885 & 19.7 & 12.3 & 0.760 & 95.6 & 0.50, 2.33\\
D1 + $G_4$           & 4,4,4,4 & 1     & 20   & --   & 1     & 100  & 0.49, 2.25 \;/\; 0.50, 2.27\\
D2 + DF              & 5,5,5,5 & 1     & 20   & --   & 1     & 100  & 0.50, $\le3$\\
\bottomrule
\end{tabular}
\end{table}

The constraints asked for in the introduction, collected:

\begin{enumerate}[label=\textbf{C\arabic*.},leftmargin=*]
\item \textbf{Balance (necessary, any rules).} Every prefix must satisfy
  \eqref{eq:balance}, and in the long run $(1-q)\nu_+\ge(1+\varepsilon)\bar\delta$,
  with a margin $\varepsilon$ for human waste and failures. Half the solutions can be
  negative only if the mean positive solution is about twice the mean number of
  operations. Example D1: at most 3 operations ($\bar\delta=\frac73$), solutions in
  $\{-6,\dots,-1\}\cup\{3,\dots,8\}$ ($\nu_+=5.5$), drift $+\frac5{12}$.
\item \textbf{Matching (with random rewards).} Every brick needs $\mu_j\ge0$. Either
  split rewards in proportion to demand (\cref{rem:matching}) or use deficit filling.
  Use an unbiased shuffle in any case (\cref{rem:shuffle}).
\item \textbf{Sure guarantee, option 1.} Deficit filling with the Lindley condition
  \eqref{eq:main}. Concretely: every negative solution is followed by a positive one
  repaying both, $\delta\le3$, $\abs{k_b},\abs k\le2$ and $c_0=(5,5,5,5)$. This keeps
  the original solutions $\pm\{1,\dots,6\}$ (\cref{cor:alternating}).
\item \textbf{Sure guarantee, option 2.} Keep the random split for its fun, and let the
  generator see the stock ($G_a$, \cref{thm:adaptive}) on top of a generator that already
  satisfies C1. The distortion is then barely measurable.
\item \textbf{Give negative bombs a purpose.} Under any sure guarantee they clear only
  the trash caused by human mistakes (\cref{prop:paradox}). Accept this, or give
  negative solutions mostly through $\pm$-bombs (\cref{thm:stabilise}).
\item \textbf{Extensions.}
  \begin{itemize}[nosep]
  \item $\dbl$ makes order matter and yields bombs that are multiples of $2^m$.
  \item $\hlv$ makes large constants cheap and the cheapest demand non-unique, so the
    player has to choose which bricks to spend.
  \item $\sq$ gives self-financing $\pm$-bombs; about a third of quadratics would already
    cancel the current deficit.
  \end{itemize}
  Each new brick adds its own budget line ($\E r(j)\ge\E d(j)$).
\end{enumerate}

\appendix
\section{Methods}\label{app:methods}

\paragraph{Exact computations.}
\begin{itemize}[nosep]
\item The joint law of $(k_b,k)$ after rejection, the demand moments of
  \cref{tab:demand} and the reward laws of the split rule were computed with rational
  arithmetic. For D0, each pair with $\abs{k_b}=1$ and $\abs k\ge1$ has probability $\frac5{101}$;
  each pair with $\abs{k_b}=2$ (then $\abs k\le2$), or with $k_b=0$ and $\abs k\in\{2,3\}$,
  has $\frac5{202}$; the cost-one pairs have $\frac1{101}$ ($k_b=0$) or $\frac2{101}$
  ($k=0$).
\item $\bar p_N$ in \cref{thm:allocfree} is an exact forward recursion over the total stock.
\item The adjustment coefficients solve the stated moment equations by bisection.
\end{itemize}

\paragraph{Exhaustive search.} Breadth-first search from the solved state $(1,0,1)$ over
the inverse operations, with exact rationals, found $17\,474$ states within 9 operations,
separately for $\{\pm\mathsf x,\pm\mathsf 1,\dbl\}$ and for
$\{\pm\mathsf x,\pm\mathsf 1,\hlv\}$. On every state the distance equals the formula of
\cref{thm:double} or \cref{thm:halve}. With both scalings, for $\Delta=1$, the distance
equals $\abs u+\abs\beta$ whenever this is at most 5, and the first detour gain appears
at $\abs u+\abs\beta=6$ (\cref{prop:detour}).

\paragraph{Simulation.} The simulated player is perfect in the sense of \cref{def:play}.
An unsolvable equation becomes trash, and play continues with the next one.
\begin{itemize}[nosep]
\item Rounds were generated exactly as in the game's code (D0) or as described for D1/D2.
\item The split rule is the game's, with a Fisher--Yates shuffle except where the biased
  sort is stated.
\item Trash stacking used the measured block widths (15{,}000 equations rendered in
  headless Chrome at five screen sizes) and uniformly random columns.
\item The screen sizes and the shuffle bias were measured on the actual page and in
  Node~22 (V8).
\end{itemize}

\end{document}
